\section{Limitations and Threats to Validity}
\label{sec:limitations}

\paragraph{What would refute the claim.} The position fails if a component of a working harness is found that (i) satisfies none of the five definitions, (ii) is not the answer to one of the six concerns within a primitive, and (iii) is not a nested harness. We invite reports of such a component. Three classes seemed most likely to produce one before the deconstruction: learned components inside the harness, search or decoding loops, and control flow generated by the model at run time. The first and third arose in the six harnesses and were placed by the rules of \S\ref{sec:rules}; search and decoding loops did not arise and remain untested.

\paragraph{Selection.} Six harnesses, five of them coding agents and one a workflow framework used for a coding example. Harnesses for other domains (browsing, voice, robotics, customer service) were not examined, and the design spaces in \S\ref{sec:primitives} are drawn from what these six do. The selection criterion was wide use and the availability of primary sources; it is not a random sample.

\paragraph{Documentation-only analysis.} Claude Code's source is not published. Its deconstruction rests on its documentation and Anthropic's published descriptions, which state what the harness does but not always how. Negative definitions are harder to confirm from documentation than from source, and a documented feature may differ from the shipped one.

\paragraph{The harness is not always in one place.} Two of the six run part of their harness behind the provider's API (\S\ref{sec:summary}). A deconstruction from client source or documentation describes the local part; the part on the provider's side is visible only through its effects. As providers absorb more harness functions, this gap will grow.

\paragraph{Products change.} Every assignment is pinned to a commit or a documentation date. Several of the six harnesses change weekly. The deconstruction describes them on those dates, and a reader checking it later should use the pinned versions.

\paragraph{Who made the assignments.} The assignments were made by the author with an AI coding agent as a research instrument (see the disclosure), against the definitions of \S\ref{sec:primitives} and the rules of \S\ref{sec:rules}, and each cites its evidence. No second, independent decomposition was made, so the agreement between two people applying the definitions is unmeasured. The rules in \S\ref{sec:rules} are where disagreement would most likely arise: they follow from the definitions, but another reader could draw some boundaries elsewhere.

\paragraph{Stipulated boundaries.} Backup versus routing, receiving versus rendering, the reader of a record and the consumer of an answer are rules we had to state. Their value is in making each choice explicit and checkable, not in being the only possible choice.

\paragraph{No experimental test.} The paper does not vary one primitive at a time and measure the effect on performance. \citet{fan2026empirical} do this for three decisions in a coding harness; a study across all five primitives would test whether the vocabulary captures the factors that matter, and is future work. The one quantitative result we cite (\S\ref{sec:uses}, footnote) comes from the companion book and is not part of this paper's argument.

\paragraph{The model is held fixed.} Training a model on traces from a harness moves work between the two components. We note this in \S\ref{sec:uses} and do not analyze the trade-off.
